\begin{abstract}
We present a comprehensive geometric analysis of magic squares through novel computational frameworks that reveal their intrinsic mathematical structure. By embedding magic squares as point clouds in $\mathbb{R}^3$, we discover a rich geometric landscape characterized by consistent convex hull properties and elegant covariance relationships. Our analysis of all 880 unique $4 \times 4$ magic squares demonstrates consistent geometric properties: convex hull volumes ranging from 58.0 to 97.3 cubic units (mean $84.1 \pm 7.1$, standard deviation), vertex counts from 8 to 16 (mean $12.2 \pm 1.4$), and crystallization into exactly 21 distinct geometric configurations. We establish a fundamental zero-covariance theorem showing that magic squares are precisely those permutation arrangements where the covariances between spatial coordinates and entry values vanish simultaneously---a property we prove is both necessary and sufficient. Remarkably, even a single element swap destroys this geometric balance entirely. Through Earth Mover's Distance analysis, we identify natural clustering with nearest-neighbor distances of $0.021 \pm 0.011$, revealing structured geometric relationships. These findings establish magic squares as a paradigmatic example of discrete mathematical objects with rich geometric structure, opening new avenues for analyzing combinatorial problems through geometric lenses. The covariance formulation enables gradient-based generation algorithms achieving $73\times$ speedup over backtracking for $n=4$. All results are validated exclusively for $n=4$; extension to larger orders remains an open problem.
\end{abstract}
